\section{Frobenius's Theorem: the Frobenius Kernel is Normal}

Throughout, $G$ is a finite group acting faithfully and transitively on a finite
set $X$ with $|X| \ge 2$, all point stabilisers are non-trivial, and $G$
satisfies the \emph{Frobenius condition}: no non-identity element of $G$ fixes
more than one point of $X$. We fix a base point $x_0 \in X$ and write
$H = \mathrm{Stab}_G(x_0)$ for its stabiliser. Frobenius's theorem (1901) asserts
that the \emph{Frobenius kernel}
$K = \{1\} \cup \{ g \in G : g \text{ fixes no point of } X \}$
is a normal subgroup of $G$. The only known proof is character-theoretic.

\subsection{The Frobenius kernel and the partition of $G$}

\begin{definition}[Frobenius group hypotheses]
    \label{def:frobenius-hypotheses}
    \lean{LeanEval.GroupTheory.FrobeniusHypotheses}
    \leanfile{LeanEval/GroupTheory/Frobenius/Basic.lean}
    \leanok
    The \emph{Frobenius hypotheses} on a finite group $G$ acting on a finite set
    $X$ with base point $x_0$ and $H = \mathrm{Stab}_G(x_0)$ are the conjunction
    of: (i) the action is faithful; (ii) the action is transitive; (iii)
    $|X| \ge 2$; (iv) every point stabiliser is non-trivial,
    $\mathrm{Stab}_G(x) \neq \{1\}$ for all $x$; and (v) the \emph{Frobenius
    condition}: no non-identity element of $G$ fixes two distinct points of $X$.
    All subsequent statements assume this bundle of hypotheses.
\end{definition}

\begin{definition}[Frobenius kernel]
    \label{def:frobenius-kernel}
    \lean{LeanEval.GroupTheory.frobeniusKernel}
    \leanfile{LeanEval/GroupTheory/Frobenius/Basic.lean}
    \leanok
    The \emph{Frobenius kernel} of the action is the set
    $K = \{1\} \cup \{ g \in G : \forall x \in X,\ g \cdot x \neq x \}$,
    consisting of the identity together with all elements acting without a fixed
    point.
\end{definition}

\begin{lemma}[Stabiliser of a moved point]
    \label{lem:stabilizer-conj}
    For every $g \in G$ and $x \in X$, the stabiliser of $g \cdot x$ is the
    conjugate $g\, \mathrm{Stab}_G(x)\, g^{-1}$ of the stabiliser of $x$.
\end{lemma}
\begin{proof}
    This is the standard conjugation formula for stabilisers under a group
    action: $g_0 \in \mathrm{Stab}_G(g\cdot x)$ iff $g^{-1} g_0 g \in
    \mathrm{Stab}_G(x)$.
\end{proof}

\begin{lemma}[Trivial intersection of distinct conjugate stabilisers]
    \label{lem:trivial-intersection}
    \lean{LeanEval.GroupTheory.stabilizer_inf_eq_bot}
    \leanfile{LeanEval/GroupTheory/Frobenius/Basic.lean}
    \uses{lem:stabilizer-conj, def:frobenius-hypotheses}
    If $x, y \in X$ are distinct, then $\mathrm{Stab}_G(x) \cap
    \mathrm{Stab}_G(y) = \{1\}$. Equivalently, for $g \notin H$ one has
    $H \cap gHg^{-1} = \{1\}$.
\end{lemma}
\begin{proof}
    \uses{lem:stabilizer-conj}
    An element of $\mathrm{Stab}_G(x) \cap \mathrm{Stab}_G(y)$ fixes both $x$ and
    $y$; by the Frobenius condition a non-identity element cannot fix two
    distinct points, so the element is $1$. The conjugate form follows from
    Lemma~\ref{lem:stabilizer-conj}, taking $y = g \cdot x_0 \neq x_0$ for
    $g \notin H$.
\end{proof}

\begin{lemma}[Dichotomy for non-identity elements]
    \label{lem:partition}
    \lean{LeanEval.GroupTheory.partition}
    \leanfile{LeanEval/GroupTheory/Frobenius/Basic.lean}
    \uses{def:frobenius-kernel, def:frobenius-hypotheses}
    Every $g \in G$ with $g \neq 1$ either fixes no point of $X$ (so $g \in K$),
    or fixes exactly one point of $X$. In the latter case $g$ lies in a unique
    point stabiliser, namely $\mathrm{Stab}_G(x)$ for the unique fixed point $x$.
\end{lemma}
\begin{proof}
    \uses{def:frobenius-kernel}
    By the Frobenius condition a non-identity $g$ fixes at most one point. So
    either it fixes none, putting it in $K$ by
    Definition~\ref{def:frobenius-kernel}, or it fixes exactly one point $x$,
    and then $g \in \mathrm{Stab}_G(x)$, with $x$ unique.
\end{proof}

\subsection{Class functions, induction and inner products}

\begin{definition}[Class function]
    \label{def:class-function}
    \lean{LeanEval.GroupTheory.IsClassFunction}
    \leanfile{LeanEval/GroupTheory/Frobenius/Character.lean}
    \leanok
    A \emph{class function} on a finite group $\Gamma$ (valued in $\mathbb{C}$) is
    a function $\varphi : \Gamma \to \mathbb{C}$ that is constant on conjugacy
    classes: $\varphi(s\gamma s^{-1}) = \varphi(\gamma)$ for all
    $s, \gamma \in \Gamma$.
\end{definition}

\begin{definition}[Inner product of class functions]
    \label{def:inner-product}
    \lean{LeanEval.GroupTheory.classInner}
    \leanfile{LeanEval/GroupTheory/Frobenius/Character.lean}
    \leanok
    For functions $\alpha, \beta : \Gamma \to \mathbb{C}$ on a finite group
    $\Gamma$, set
    $\langle \alpha, \beta \rangle_\Gamma = \frac{1}{|\Gamma|} \sum_{\gamma \in
    \Gamma} \alpha(\gamma) \overline{\beta(\gamma)}$.
\end{definition}

\begin{lemma}[Eigenvalues of a representation element are roots of unity]
    \label{lem:rep-element-eigenvalues-roots-of-unity}
    \lean{LeanEval.GroupTheory.rep_eigenvalues_rootsOfUnity}
    \leanfile{LeanEval/GroupTheory/Frobenius/Character.lean}
    \uses{def:irreducible-character}
    Let $\rho$ be a finite-dimensional complex representation of a finite group
    $\Gamma$ with character $\chi$ and $d = \dim\rho = \chi(1)$. For every
    $\gamma \in \Gamma$ the operator $\rho(\gamma)$ is diagonalisable, and its
    eigenvalues $\lambda_1, \dots, \lambda_d$ (listed with multiplicity) are roots
    of unity; in particular each $|\lambda_i| = 1$ and
    $\lambda_i^{-1} = \overline{\lambda_i}$, while
    $\chi(\gamma) = \sum_{i} \lambda_i$.
\end{lemma}
\begin{proof}
    Since $\Gamma$ is finite, $\gamma$ has finite order $n = \mathrm{ord}(\gamma)$
    and $\rho(\gamma)^n = \rho(\gamma^n) = \rho(1) = 1$ (via \texttt{map\_pow},
    \texttt{pow\_orderOf\_eq\_one}, \texttt{map\_one}). Thus $\rho(\gamma)$ is a
    root of $X^n - 1$, a polynomial with $n$ distinct roots over $\mathbb{C}$, so
    its minimal polynomial is squarefree and $\rho(\gamma)$ is diagonalisable.
    Each eigenvalue $\lambda$ satisfies $\lambda^n = 1$, i.e.\ $\lambda$ is an
    $n$-th root of unity, whence $|\lambda| = 1$
    (\texttt{Complex.norm\_eq\_one\_of\_mem\_rootsOfUnity}) and
    $\lambda^{-1} = \overline{\lambda}$ (\texttt{Complex.inv\_eq\_conj}). The
    trace of a diagonalisable operator is the sum of its eigenvalues with
    multiplicity, so $\chi(\gamma) = \sum_i \lambda_i$.
\end{proof}

\begin{lemma}[Character of an inverse is the conjugate]
    \label{lem:char-inverse-conjugate}
    \lean{LeanEval.GroupTheory.char_inv_conj}
    \leanfile{LeanEval/GroupTheory/Frobenius/Character.lean}
    \uses{def:irreducible-character, lem:rep-element-eigenvalues-roots-of-unity}
    For the character $\chi$ of a finite-dimensional complex representation of a
    finite group $\Gamma$, $\chi(\gamma^{-1}) = \overline{\chi(\gamma)}$ for all
    $\gamma \in \Gamma$.
\end{lemma}
\begin{proof}
    \uses{lem:rep-element-eigenvalues-roots-of-unity}
    By Lemma~\ref{lem:rep-element-eigenvalues-roots-of-unity},
    $\rho(\gamma)$ has eigenvalues $\lambda_1, \dots, \lambda_d$ (roots of unity)
    with $\chi(\gamma) = \sum_i \lambda_i$. The eigenvalues of
    $\rho(\gamma^{-1}) = \rho(\gamma)^{-1}$ are the $\lambda_i^{-1}$, so
    $\chi(\gamma^{-1}) = \sum_i \lambda_i^{-1} = \sum_i \overline{\lambda_i} =
    \overline{\sum_i \lambda_i} = \overline{\chi(\gamma)}$, using
    $\lambda_i^{-1} = \overline{\lambda_i}$ for roots of unity.
\end{proof}

\begin{lemma}[Inner product as a character sum]
    \label{lem:inner-product-as-character-sum}
    \lean{LeanEval.GroupTheory.inner_as_char_sum}
    \leanfile{LeanEval/GroupTheory/Frobenius/Character.lean}
    \uses{def:inner-product, lem:char-inverse-conjugate, def:irreducible-character}
    For characters $\chi, \psi$ of a finite group $\Gamma$,
    $\langle \chi, \psi \rangle_\Gamma = \frac{1}{|\Gamma|}\sum_{\gamma \in
    \Gamma} \chi(\gamma)\,\psi(\gamma^{-1})$. In particular it matches the
    Mathlib sum $\sum_{\gamma} \chi(\gamma)\psi(\gamma^{-1})$ used in
    \texttt{char\_orthonormal} and \texttt{simple\_iff\_char\_is\_norm\_one}.
\end{lemma}
\begin{proof}
    \uses{lem:char-inverse-conjugate}
    Substitute $\overline{\psi(\gamma)} = \psi(\gamma^{-1})$ from
    Lemma~\ref{lem:char-inverse-conjugate} into
    Definition~\ref{def:inner-product}.
\end{proof}

\begin{lemma}[Inner product equals dimension of equivariant maps]
    \label{lem:inner-product-eq-dim-hom}
    \lean{LeanEval.GroupTheory.inner_eq_dim_hom}
    \leanfile{LeanEval/GroupTheory/Frobenius/Character.lean}
    \uses{lem:inner-product-as-character-sum, def:irreducible-character}
    For representations $V, W$ of a finite group $\Gamma$ over $\mathbb{C}$ with
    characters $\chi_V, \chi_W$,
    $\langle \chi_V, \chi_W \rangle_\Gamma = \dim_{\mathbb{C}}
    \mathrm{Hom}_\Gamma(W, V)$, the dimension of the space of $\Gamma$-equivariant
    maps $W \to V$.
\end{lemma}
\begin{proof}
    \uses{lem:inner-product-as-character-sum}
    By Lemma~\ref{lem:inner-product-as-character-sum} the inner product equals
    $\frac{1}{|\Gamma|}\sum_\gamma \chi_V(\gamma)\chi_W(\gamma^{-1})$, which is
    exactly Mathlib's \texttt{Representation.card\_inv\_mul\_sum\_char\_mul\_char\_eq\_finrank}
    (equivalently \texttt{FDRep.scalar\_product\_char\_eq\_finrank\_equivariant}),
    equal to $\dim_{\mathbb{C}}\mathrm{Hom}_\Gamma(W,V)$.
\end{proof}

\begin{definition}[Induced class function]
    \label{def:induced-class-function}
    \lean{LeanEval.GroupTheory.indClassFun}
    \leanfile{LeanEval/GroupTheory/Frobenius/Induction.lean}
    \uses{def:class-function}
    \leanok
    For a class function $\varphi$ on $H \le G$, the \emph{induced class
    function} $\mathrm{Ind}_H^G \varphi : G \to \mathbb{C}$ is
    $(\mathrm{Ind}_H^G \varphi)(g) = \frac{1}{|H|} \sum_{\substack{x \in G \\
    x^{-1} g x \in H}} \varphi(x^{-1} g x)$,
    summing only over the conjugators carrying $g$ into $H$.
\end{definition}

\begin{lemma}[Coset basis of the induced representation]
    \label{lem:ind-coset-basis}
    \lean{LeanEval.GroupTheory.ind_coset_basis}
    \leanfile{LeanEval/GroupTheory/Frobenius/Induction.lean}
    \uses{def:irreducible-character}
    Let $\iota : H \hookrightarrow G$ be the inclusion and $V$ a
    finite-dimensional $H$-representation. Fix a left transversal $T$ of $H$ in
    $G$ (\texttt{Subgroup.LeftTransversal}, which exists by
    \texttt{Subgroup.instInhabitedLeftTransversal}) and a basis $(e_j)$ of $V$.
    Then the induced representation
    $\mathrm{ind}_\iota V = (k[G] \otimes_k V)_H$ (Mathlib's
    \texttt{Representation.ind}) has $k$-basis $\{\,\llbracket t \otimes e_j
    \rrbracket : t \in T,\ j\,\}$, and for $g \in G$ the action satisfies
    $g \cdot \llbracket t \otimes v \rrbracket = \llbracket (gt) \otimes v
    \rrbracket$; writing $gt = t' h$ with $t' \in T$ the representative of the
    coset $gtH$ and $h = (t')^{-1} g t \in H$, this equals
    $\llbracket t' \otimes \rho_V(h) v \rrbracket$.
\end{lemma}
\begin{proof}
    \uses{def:irreducible-character}
    By \texttt{Representation.ind\_apply} the $G$-action on
    $(k[G] \otimes_k V)_H$ sends $\llbracket t \otimes v \rrbracket$ to
    $\llbracket (gt) \otimes v \rrbracket$. Each coset of $H$ in $G$ has a unique
    representative in the transversal $T$, so the classes $\llbracket t \otimes
    e_j \rrbracket$ are a $k$-basis. Rewriting $gt = t' h$ with $t' \in T$ and
    $h \in H$ and using $H$-invariance of the coinvariants gives
    $\llbracket (gt) \otimes v \rrbracket = \llbracket t' \otimes \rho_V(h) v
    \rrbracket$.
\end{proof}

\begin{lemma}[Block-monomial matrix of the induced action]
    \label{lem:ind-block-matrix}
    \lean{LeanEval.GroupTheory.ind_block_matrix}
    \leanfile{LeanEval/GroupTheory/Frobenius/Induction.lean}
    \uses{lem:ind-coset-basis}
    In the coset basis of Lemma~\ref{lem:ind-coset-basis}, the matrix of
    $\rho_{\mathrm{ind}}(g)$ is block-monomial: the block sending the
    $t$-component to the $t'$-component is non-zero only when $t' H = g t H$. In
    particular the diagonal block at $t \in T$ is non-zero iff $g$ fixes the coset
    $tH$, equivalently $t^{-1} g t \in H$, in which case that diagonal block is
    $\rho_V(t^{-1} g t)$.
\end{lemma}
\begin{proof}
    \uses{lem:ind-coset-basis}
    By Lemma~\ref{lem:ind-coset-basis}, $g$ carries the block at $t$ into the
    block at $t'$, where $t'$ is the representative of $gtH$; hence the only
    non-zero block out of $t$ lands at $t' H = gtH$. The diagonal block ($t' = t$)
    occurs exactly when $gtH = tH$, i.e.\ $t^{-1} g t \in H$, and there the action
    is $v \mapsto \rho_V(t^{-1} g t) v$.
\end{proof}

\begin{lemma}[Trace of the induced action as a sum over conjugators]
    \label{lem:ind-trace-sum}
    \lean{LeanEval.GroupTheory.ind_trace_sum}
    \leanfile{LeanEval/GroupTheory/Frobenius/Induction.lean}
    \uses{lem:ind-block-matrix, def:induced-class-function}
    With $\varphi = \chi_V$ the character of $V$, the character of
    $\mathrm{ind}_\iota V$ satisfies, for every $g \in G$,
    $\chi_{\mathrm{ind}_\iota V}(g) = \sum_{\substack{t \in T \\ t^{-1} g t \in H}}
    \varphi(t^{-1} g t) = \frac{1}{|H|} \sum_{\substack{x \in G \\ x^{-1} g x \in
    H}} \varphi(x^{-1} g x)$.
\end{lemma}
\begin{proof}
    \uses{lem:ind-block-matrix, def:induced-class-function}
    The trace of the block-monomial matrix of
    Lemma~\ref{lem:ind-block-matrix} is the sum of the traces of its diagonal
    blocks, and only the cosets fixed by $g$ contribute; each such block has trace
    $\mathrm{tr}\,\rho_V(t^{-1} g t) = \varphi(t^{-1} g t)$. This gives the sum
    over the transversal. Passing from $T$ to all of $G$: writing $x = t h$ with
    $h \in H$, $\varphi((th)^{-1} g (th)) = \varphi(h^{-1}(t^{-1} g t) h) =
    \varphi(t^{-1} g t)$ since $\varphi$ is a class function on $H$, and
    $t^{-1} g t \in H \iff x^{-1} g x \in H$; each transversal term is thus
    repeated $|H|$ times, giving the $\frac{1}{|H|}$-normalised sum over $G$.
\end{proof}

\begin{lemma}[Induced class function is the character of $\mathrm{ind}$]
    \label{lem:induced-is-char-of-ind}
    \lean{LeanEval.GroupTheory.induced_is_char_of_ind}
    \leanfile{LeanEval/GroupTheory/Frobenius/Induction.lean}
    \uses{def:induced-class-function, def:irreducible-character, lem:ind-trace-sum}
    Let $\iota : H \hookrightarrow G$ be the inclusion. If $\varphi$ is the
    character of an $H$-representation $V$, then $\mathrm{Ind}_H^G \varphi$ is the
    character of the induced $G$-representation $\mathrm{ind}_\iota V$ (Mathlib's
    \texttt{Representation.ind}). Consequently induction of characters is realised
    by \texttt{Rep.ind} at the level of representations.
\end{lemma}
\begin{proof}
    \uses{lem:ind-trace-sum}
    By Lemma~\ref{lem:ind-trace-sum} the character of $\mathrm{ind}_\iota V$ at
    $g$ equals $\frac{1}{|H|}\sum_{x^{-1} g x \in H}\varphi(x^{-1} g x)$, which is
    exactly $(\mathrm{Ind}_H^G \varphi)(g)$ of
    Definition~\ref{def:induced-class-function}.
\end{proof}

\begin{lemma}[Induction yields a class function]
    \label{lem:induced-is-class-function}
    \lean{LeanEval.GroupTheory.indClassFun_isClassFunction}
    \leanfile{LeanEval/GroupTheory/Frobenius/Induction.lean}
    \uses{def:induced-class-function, def:class-function}
    For any class function $\varphi$ on $H$, $\mathrm{Ind}_H^G \varphi$ is a
    class function on $G$.
\end{lemma}
\begin{proof}
    \uses{def:induced-class-function}
    Argue directly by reindexing the defining sum, without appeal to the
    completeness of characters. Fix $s, g \in G$. Then
    $(\mathrm{Ind}_H^G\varphi)(s g s^{-1}) = \frac{1}{|H|}\sum_{x :\, x^{-1}(s g
    s^{-1})x \in H}\varphi(x^{-1} s g s^{-1} x)$. Substituting $x = s y$ (a
    bijection of $G$) gives $x^{-1} s g s^{-1} x = y^{-1} g y$, so the index set
    and summand match those of $(\mathrm{Ind}_H^G\varphi)(g)$ term-by-term under
    $x \leftrightarrow y$. Hence $(\mathrm{Ind}_H^G\varphi)(s g s^{-1}) =
    (\mathrm{Ind}_H^G\varphi)(g)$, i.e.\ $\mathrm{Ind}_H^G\varphi$ is constant on
    conjugacy classes.
\end{proof}

\begin{lemma}[Induced value at the identity]
    \label{lem:induced-value-at-one}
    \lean{LeanEval.GroupTheory.indClassFun_one}
    \leanfile{LeanEval/GroupTheory/Frobenius/Induction.lean}
    \uses{def:induced-class-function}
    For any class function $\varphi$ on $H$, $(\mathrm{Ind}_H^G \varphi)(1) =
    [G:H]\,\varphi(1)$.
\end{lemma}
\begin{proof}
    \uses{def:induced-class-function}
    At $g = 1$ every conjugate $x^{-1}\cdot 1\cdot x = 1 \in H$, so the defining
    sum runs over all $x \in G$ and equals
    $\frac{1}{|H|}\cdot |G|\cdot \varphi(1) = [G:H]\,\varphi(1)$.
\end{proof}

\begin{lemma}[Induced inner product as a $\mathrm{Hom}$-dimension]
    \label{lem:ind-inner-as-hom-dim}
    \lean{LeanEval.GroupTheory.ind_inner_as_hom_dim}
    \leanfile{LeanEval/GroupTheory/Frobenius/Induction.lean}
    \uses{lem:induced-is-char-of-ind, lem:inner-product-eq-dim-hom}
    Let $V, W$ realise the characters $\varphi$ of $H$ and $\psi$ of $G$. Then
    $\langle \mathrm{Ind}_H^G\varphi, \psi\rangle_G = \dim_{\mathbb{C}}
    \mathrm{Hom}_G(W, \mathrm{ind}_\iota V)$.
\end{lemma}
\begin{proof}
    \uses{lem:induced-is-char-of-ind, lem:inner-product-eq-dim-hom}
    By Lemma~\ref{lem:induced-is-char-of-ind}, $\mathrm{Ind}_H^G\varphi$ is the
    character of $\mathrm{ind}_\iota V$. Apply
    Lemma~\ref{lem:inner-product-eq-dim-hom} to the $G$-representations
    $\mathrm{ind}_\iota V$ and $W$.
\end{proof}

\begin{lemma}[Restricted inner product as a $\mathrm{Hom}$-dimension]
    \label{lem:res-inner-as-hom-dim}
    \lean{LeanEval.GroupTheory.res_inner_as_hom_dim}
    \leanfile{LeanEval/GroupTheory/Frobenius/Induction.lean}
    \uses{lem:inner-product-eq-dim-hom}
    Let $V, W$ realise the characters $\varphi$ of $H$ and $\psi$ of $G$. Then
    $\langle \varphi, \psi|_H\rangle_H = \dim_{\mathbb{C}}
    \mathrm{Hom}_H(\mathrm{res}_\iota W, V)$, where $\mathrm{res}_\iota W$ is $W$
    viewed as an $H$-representation.
\end{lemma}
\begin{proof}
    \uses{lem:inner-product-eq-dim-hom}
    The character of $\mathrm{res}_\iota W$ is $\psi|_H$. Apply
    Lemma~\ref{lem:inner-product-eq-dim-hom} to the $H$-representations $V$ and
    $\mathrm{res}_\iota W$.
\end{proof}

\begin{lemma}[Frobenius reciprocity for characters]
    \label{lem:frobenius-reciprocity-char}
    \lean{LeanEval.GroupTheory.frobenius_reciprocity_char}
    \leanfile{LeanEval/GroupTheory/Frobenius/Induction.lean}
    \uses{lem:ind-inner-as-hom-dim, lem:res-inner-as-hom-dim}
    For a character $\varphi$ of $H$ and a character $\psi$ of $G$,
    $\langle \mathrm{Ind}_H^G \varphi, \psi \rangle_G =
    \langle \varphi, \psi|_H \rangle_H$.
\end{lemma}
\begin{proof}
    \uses{lem:ind-inner-as-hom-dim, lem:res-inner-as-hom-dim}
    Let $V, W$ realise $\varphi, \psi$. By
    Lemma~\ref{lem:ind-inner-as-hom-dim} the left side is $\dim
    \mathrm{Hom}_G(W, \mathrm{ind}_\iota V)$ and by
    Lemma~\ref{lem:res-inner-as-hom-dim} the right side is $\dim
    \mathrm{Hom}_H(\mathrm{res}_\iota W, V)$. The Frobenius reciprocity
    equivalence \texttt{Rep.indResHomEquiv} gives a linear isomorphism
    $\mathrm{Hom}_G(\mathrm{ind}_\iota V, W) \cong \mathrm{Hom}_H(V,
    \mathrm{res}_\iota W)$; symmetry of $\dim\mathrm{Hom}$ for these
    finite-dimensional spaces equates the two dimensions.
\end{proof}

\begin{lemma}[Frobenius reciprocity]
    \label{lem:frobenius-reciprocity}
    \lean{LeanEval.GroupTheory.frobenius_reciprocity}
    \leanfile{LeanEval/GroupTheory/Frobenius/Induction.lean}
    \uses{def:induced-class-function, def:inner-product, lem:frobenius-reciprocity-char, lem:irr-characters-orthonormal-basis}
    For a class function $\varphi$ on $H$ and a class function $\psi$ on $G$,
    $\langle \mathrm{Ind}_H^G \varphi, \psi \rangle_G =
    \langle \varphi, \psi|_H \rangle_H$,
    where $\psi|_H$ is the restriction of $\psi$ to $H$.
\end{lemma}
\begin{proof}
    \uses{lem:frobenius-reciprocity-char, lem:irr-characters-orthonormal-basis}
    Both sides are sesquilinear in $(\varphi, \psi)$ and induction and
    restriction are linear, so the identity extends from the case of genuine
    characters (Lemma~\ref{lem:frobenius-reciprocity-char}) to arbitrary class
    functions. Writing each class function as a $\mathbb{C}$-combination of
    irreducible characters \emph{relies on the completeness of the irreducible
    characters}, i.e.\ on the accepted assumption
    Lemma~\ref{lem:irr-characters-orthonormal-basis}; this dependence is made
    explicit here.
\end{proof}

\begin{lemma}[Induced function vanishes on the kernel]
    \label{lem:induced-vanishing-on-kernel}
    \lean{LeanEval.GroupTheory.indClassFun_vanishing_on_kernel}
    \leanfile{LeanEval/GroupTheory/Frobenius/Induction.lean}
    \uses{def:induced-class-function, lem:partition, def:frobenius-kernel}
    Let $\varphi$ be any class function on $H$. Then $(\mathrm{Ind}_H^G
    \varphi)(g) = 0$ for every fixed-point-free $g$ (every $g \in K \setminus
    \{1\}$).
\end{lemma}
\begin{proof}
    \uses{lem:partition, def:frobenius-kernel}
    A fixed-point-free $g$ has no conjugate lying in $H$: any $x^{-1} g x \in H$
    would fix $x_0$, so $g$ would fix $x \cdot x_0$, contradicting
    Lemma~\ref{lem:partition}. The defining sum is therefore empty and the value
    is $0$.
\end{proof}

\begin{lemma}[Conjugators carrying a non-identity element of $H$ into $H$ lie in $H$]
    \label{lem:conjugators-into-H-are-H}
    \lean{LeanEval.GroupTheory.conjugators_into_H}
    \leanfile{LeanEval/GroupTheory/Frobenius/Basic.lean}
    \uses{lem:trivial-intersection, def:frobenius-hypotheses}
    Let $h \in H$ with $h \neq 1$. If $x \in G$ satisfies $x^{-1} h x \in H$, then
    $x \in H$. Consequently $\{x \in G : x^{-1} h x \in H\} = H$.
\end{lemma}
\begin{proof}
    \uses{lem:trivial-intersection}
    Suppose $x^{-1} h x \in H = \mathrm{Stab}_G(x_0)$. Then $h$ fixes $x \cdot
    x_0$, and $h \in H$ fixes $x_0$, so $h \in \mathrm{Stab}_G(x_0) \cap
    \mathrm{Stab}_G(x \cdot x_0)$. If $x \cdot x_0 \neq x_0$ this intersection is
    trivial by Lemma~\ref{lem:trivial-intersection}, contradicting $h \neq 1$.
    Hence $x \cdot x_0 = x_0$, i.e.\ $x \in H$. The reverse inclusion is immediate
    since $H$ is a subgroup, giving the stated equality.
\end{proof}

\begin{lemma}[Induced function restricts to $H$ when $\varphi(1)=0$]
    \label{lem:induced-restriction-to-H}
    \lean{LeanEval.GroupTheory.indClassFun_restriction_to_H}
    \leanfile{LeanEval/GroupTheory/Frobenius/Induction.lean}
    \uses{def:induced-class-function, lem:conjugators-into-H-are-H}
    Let $\varphi$ be a class function on $H$ with $\varphi(1) = 0$. Then
    $(\mathrm{Ind}_H^G \varphi)(h) = \varphi(h)$ for every $h \in H$.
\end{lemma}
\begin{proof}
    \uses{lem:conjugators-into-H-are-H}
    If $h = 1$ both sides are $\varphi(1) = 0$. For $h \neq 1$, by
    Lemma~\ref{lem:conjugators-into-H-are-H} the conjugators $x$ with $x^{-1} h x
    \in H$ are exactly those in $H$, so the defining sum collapses to
    $\frac{1}{|H|}\sum_{x \in H}\varphi(x^{-1}hx) = \varphi(h)$, since $\varphi$ is
    a class function on $H$ and $|H|$ terms are each equal to $\varphi(h)$.
\end{proof}

\begin{lemma}[Induced value of a vanishing class function]
    \label{lem:induced-vanishing-restriction}
    \lean{LeanEval.GroupTheory.indClassFun_vanishing_restriction}
    \leanfile{LeanEval/GroupTheory/Frobenius/Induction.lean}
    \uses{lem:induced-vanishing-on-kernel, lem:induced-restriction-to-H}
    Let $\varphi$ be a class function on $H$ with $\varphi(1) = 0$. Then
    $(\mathrm{Ind}_H^G \varphi)(h) = \varphi(h)$ for every $h \in H$, and
    $(\mathrm{Ind}_H^G \varphi)(g) = 0$ for every $g \in K \setminus \{1\}$.
\end{lemma}
\begin{proof}
    \uses{lem:induced-vanishing-on-kernel, lem:induced-restriction-to-H}
    Combine Lemma~\ref{lem:induced-restriction-to-H} (value on $H$) and
    Lemma~\ref{lem:induced-vanishing-on-kernel} (vanishing on fixed-point-free
    elements).
\end{proof}

\subsection{Lifting an irreducible character of $H$}

\begin{definition}[Irreducible character]
    \label{def:irreducible-character}
    \lean{LeanEval.GroupTheory.IsIrreducibleCharacter}
    \leanfile{LeanEval/GroupTheory/Frobenius/Character.lean}
    \leanok
    For a finite group $\Gamma$, an \emph{irreducible character} is the character
    $\chi(\gamma) = \mathrm{tr}\,\rho(\gamma)$ of a simple (irreducible)
    finite-dimensional complex representation $\rho$ of $\Gamma$. Its degree is
    $\chi(1) = \dim \rho > 0$.
\end{definition}

\begin{lemma}[Character inner products are integers]
    \label{lem:character-inner-product-integer}
    \lean{LeanEval.GroupTheory.char_inner_product_integer}
    \leanfile{LeanEval/GroupTheory/Frobenius/Character.lean}
    \uses{def:irreducible-character, def:inner-product, lem:inner-product-as-character-sum}
    For irreducible characters $\chi_i, \chi_j$ of a finite group $\Gamma$,
    $\langle \chi_i, \chi_j \rangle_\Gamma = \delta_{ij} \in \{0,1\}$. Consequently
    $\langle \alpha, \chi \rangle_\Gamma \in \mathbb{Z}$ whenever $\alpha$ is an
    integer combination of irreducible characters and $\chi$ is irreducible.
\end{lemma}
\begin{proof}
    \uses{lem:inner-product-as-character-sum}
    By Lemma~\ref{lem:inner-product-as-character-sum} the inner product equals the
    Mathlib character sum, so orthonormality is exactly
    \texttt{FDRep.char\_orthonormal}: it is $1$ when the simple representations are
    isomorphic and $0$ otherwise. Linearity of the inner product in the first
    argument gives the integrality statement.
\end{proof}

\begin{lemma}[Irreducible characters are an orthonormal basis of class functions (accepted assumption)]
    \label{lem:irr-characters-orthonormal-basis}
    \lean{LeanEval.GroupTheory.irr_characters_orthonormal_basis}
    \leanfile{LeanEval/GroupTheory/Frobenius/Character.lean}
    \uses{def:class-function, lem:character-inner-product-integer, def:inner-product}
    \emph{This statement is taken as an explicit accepted assumption: it is
    \textbf{not} available in Mathlib.} The irreducible characters of a finite
    group $\Gamma$ form an orthonormal basis of the space of class functions on
    $\Gamma$ with respect to $\langle\cdot,\cdot\rangle_\Gamma$. In particular
    every class function is a unique $\mathbb{C}$-linear combination of irreducible
    characters, and a class function $\alpha$ with $\langle \alpha,
    \chi\rangle_\Gamma = 0$ for every irreducible $\chi$ is identically zero.
\end{lemma}
\begin{proof}
    \uses{lem:character-inner-product-integer}
    Orthonormality of the irreducible characters is
    Lemma~\ref{lem:character-inner-product-integer}
    (\texttt{FDRep.char\_orthonormal}), and this part \emph{is} in Mathlib. The
    completeness/spanning half — that the number of irreducible characters equals
    the number of conjugacy classes, i.e.\ the dimension of the space of class
    functions, so that the orthonormal set is a full basis — is \textbf{not}
    formalised in Mathlib (only \texttt{FDRep.char\_orthonormal} and
    \texttt{Representation.char\_conj} exist, not the counting theorem). We
    therefore accept this lemma as an assumption rather than as a derived result;
    every downstream lemma that invokes the basis property does so through this
    explicit assumption.
\end{proof}

\begin{definition}[Lifted class function]
    \label{def:lifted-class-function}
    \lean{LeanEval.GroupTheory.liftClassFun}
    \leanfile{LeanEval/GroupTheory/Frobenius/Induction.lean}
    \uses{def:induced-class-function, def:irreducible-character}
    \leanok
    For an irreducible character $\chi$ of $H$, put
    $\theta = \chi - \chi(1)\,\mathbf{1}_H$ (so $\theta(1) = 0$) and define the
    \emph{lift}
    $\widetilde{\chi} = \mathrm{Ind}_H^G \theta + \chi(1)\,\mathbf{1}_G$,
    a class function on $G$.
\end{definition}

\begin{lemma}[The lift is a virtual character]
    \label{lem:lift-is-virtual-character}
    \lean{LeanEval.GroupTheory.liftClassFun_isVirtualCharacter}
    \leanfile{LeanEval/GroupTheory/Frobenius/Induction.lean}
    \uses{def:lifted-class-function, lem:frobenius-reciprocity, lem:character-inner-product-integer, lem:irr-characters-orthonormal-basis}
    For each irreducible character $\chi$ of $H$, the lift $\widetilde{\chi}$ is
    an integer combination of the irreducible characters of $G$ (a virtual
    character).
\end{lemma}
\begin{proof}
    \uses{lem:frobenius-reciprocity, lem:character-inner-product-integer, lem:irr-characters-orthonormal-basis}
    By Lemma~\ref{lem:irr-characters-orthonormal-basis} the class function
    $\widetilde{\chi}$ expands as $\sum_\zeta \langle \widetilde{\chi},
    \zeta\rangle_G\,\zeta$ over the irreducible characters $\zeta$ of $G$, so it
    suffices that every coefficient
    $\langle \widetilde{\chi}, \zeta \rangle_G$ ($\zeta$ irreducible on $G$) is an
    integer. By Lemma~\ref{lem:frobenius-reciprocity},
    $\langle \mathrm{Ind}_H^G \theta, \zeta \rangle_G =
    \langle \theta, \zeta|_H \rangle_H$, an inner product of (virtual) characters
    of $H$, hence an integer by
    Lemma~\ref{lem:character-inner-product-integer}; the $\chi(1)\mathbf 1_G$ term
    contributes an integer as well.
\end{proof}

\begin{lemma}[The lift restricts to $\chi$]
    \label{lem:lift-restricts}
    \lean{LeanEval.GroupTheory.liftClassFun_restricts}
    \leanfile{LeanEval/GroupTheory/Frobenius/Induction.lean}
    \uses{def:lifted-class-function, lem:induced-vanishing-restriction}
    For every $h \in H$, $\widetilde{\chi}(h) = \chi(h)$. In particular
    $\widetilde{\chi}(1) = \chi(1)$.
\end{lemma}
\begin{proof}
    \uses{lem:induced-vanishing-restriction}
    Apply Lemma~\ref{lem:induced-vanishing-restriction} to $\theta = \chi -
    \chi(1)\mathbf 1_H$ (which satisfies $\theta(1)=0$): for $h \in H$,
    $(\mathrm{Ind}_H^G\theta)(h) = \theta(h) = \chi(h) - \chi(1)$, so
    $\widetilde{\chi}(h) = \chi(h) - \chi(1) + \chi(1) = \chi(h)$.
\end{proof}

\begin{lemma}[Self inner product of the difference $\theta$ on $H$]
    \label{lem:theta-inner-self}
    \lean{LeanEval.GroupTheory.theta_inner_self}
    \leanfile{LeanEval/GroupTheory/Frobenius/Induction.lean}
    \uses{def:lifted-class-function, lem:character-inner-product-integer}
    With $\theta = \chi - \chi(1)\mathbf 1_H$ for $\chi$ a non-trivial irreducible
    character of $H$, $\langle \theta, \theta\rangle_H = 1 + \chi(1)^2$.
\end{lemma}
\begin{proof}
    \uses{lem:character-inner-product-integer}
    Expand by sesquilinearity:
    $\langle\theta,\theta\rangle_H = \langle\chi,\chi\rangle_H -
    2\chi(1)\,\mathrm{Re}\langle\chi,\mathbf 1_H\rangle_H + \chi(1)^2\langle\mathbf
    1_H,\mathbf 1_H\rangle_H$. By orthonormality
    (Lemma~\ref{lem:character-inner-product-integer}),
    $\langle\chi,\chi\rangle_H = 1$, $\langle\mathbf 1_H,\mathbf 1_H\rangle_H = 1$,
    and $\langle\chi,\mathbf 1_H\rangle_H = 0$ as $\chi$ is non-trivial
    irreducible. Hence $\langle\theta,\theta\rangle_H = 1 + \chi(1)^2$.
\end{proof}

\begin{lemma}[Self inner product of the induced part]
    \label{lem:lift-inner-self}
    \lean{LeanEval.GroupTheory.lift_inner_self}
    \leanfile{LeanEval/GroupTheory/Frobenius/Induction.lean}
    \uses{def:lifted-class-function, lem:frobenius-reciprocity, lem:induced-restriction-to-H, lem:theta-inner-self}
    With $\theta = \chi - \chi(1)\mathbf 1_H$ for $\chi$ a non-trivial irreducible
    character of $H$, $\langle \mathrm{Ind}_H^G\theta, \mathrm{Ind}_H^G\theta
    \rangle_G = 1 + \chi(1)^2$.
\end{lemma}
\begin{proof}
    \uses{lem:frobenius-reciprocity, lem:induced-restriction-to-H, lem:theta-inner-self}
    Frobenius reciprocity (Lemma~\ref{lem:frobenius-reciprocity}) gives $\langle
    \mathrm{Ind}_H^G\theta, \mathrm{Ind}_H^G\theta\rangle_G = \langle \theta,
    (\mathrm{Ind}_H^G\theta)|_H\rangle_H$, and $(\mathrm{Ind}_H^G\theta)|_H =
    \theta$ by Lemma~\ref{lem:induced-restriction-to-H} (as $\theta(1)=0$). Hence
    it equals $\langle\theta,\theta\rangle_H = 1 + \chi(1)^2$ by
    Lemma~\ref{lem:theta-inner-self}.
\end{proof}

\begin{lemma}[Inner product of the induced part with the trivial character]
    \label{lem:lift-inner-trivial}
    \lean{LeanEval.GroupTheory.lift_inner_trivial}
    \leanfile{LeanEval/GroupTheory/Frobenius/Induction.lean}
    \uses{def:lifted-class-function, lem:frobenius-reciprocity, lem:character-inner-product-integer}
    With $\theta = \chi - \chi(1)\mathbf 1_H$ for $\chi$ a non-trivial irreducible
    character of $H$, $\langle \mathrm{Ind}_H^G\theta, \mathbf 1_G\rangle_G =
    -\chi(1)$.
\end{lemma}
\begin{proof}
    \uses{lem:frobenius-reciprocity, lem:character-inner-product-integer}
    Frobenius reciprocity gives $\langle \mathrm{Ind}_H^G\theta, \mathbf
    1_G\rangle_G = \langle \theta, \mathbf 1_G|_H\rangle_H = \langle \theta,
    \mathbf 1_H\rangle_H$. With $\theta = \chi - \chi(1)\mathbf 1_H$ and
    $\langle\chi,\mathbf 1_H\rangle_H = 0$,
    $\langle\mathbf 1_H,\mathbf 1_H\rangle_H = 1$
    (Lemma~\ref{lem:character-inner-product-integer}), this is $-\chi(1)$.
\end{proof}

\begin{lemma}[The lift has norm one]
    \label{lem:lift-norm-one}
    \lean{LeanEval.GroupTheory.lift_norm_one}
    \leanfile{LeanEval/GroupTheory/Frobenius/Induction.lean}
    \uses{def:lifted-class-function, lem:lift-inner-self, lem:lift-inner-trivial, lem:character-inner-product-integer}
    For each non-trivial irreducible character $\chi$ of $H$, $\langle
    \widetilde{\chi}, \widetilde{\chi} \rangle_G = 1$.
\end{lemma}
\begin{proof}
    \uses{lem:lift-inner-self, lem:lift-inner-trivial, lem:character-inner-product-integer}
    Expand $\widetilde{\chi} = \mathrm{Ind}_H^G\theta + \chi(1)\mathbf 1_G$ by
    sesquilinearity:
    $\langle\widetilde\chi,\widetilde\chi\rangle_G = \langle
    \mathrm{Ind}_H^G\theta,\mathrm{Ind}_H^G\theta\rangle_G + 2\chi(1)\,\langle
    \mathrm{Ind}_H^G\theta,\mathbf 1_G\rangle_G + \chi(1)^2\langle\mathbf
    1_G,\mathbf 1_G\rangle_G$. Substituting $1+\chi(1)^2$
    (Lemma~\ref{lem:lift-inner-self}), $-\chi(1)$
    (Lemma~\ref{lem:lift-inner-trivial}) and $\langle\mathbf 1_G,\mathbf
    1_G\rangle_G = 1$ yields $1 + \chi(1)^2 - 2\chi(1)^2 + \chi(1)^2 = 1$.
\end{proof}

\begin{lemma}[A norm-one virtual character is $\pm$ an irreducible character]
    \label{lem:norm-one-virtual-is-plus-minus-irr}
    \lean{LeanEval.GroupTheory.norm_one_virtual_is_plus_minus_irr}
    \leanfile{LeanEval/GroupTheory/Frobenius/Character.lean}
    \uses{def:irreducible-character, lem:character-inner-product-integer, lem:irr-characters-orthonormal-basis}
    Let $\alpha$ be a virtual character of a finite group $\Gamma$ with
    $\langle\alpha,\alpha\rangle_\Gamma = 1$. Then $\alpha = \pm\zeta$ for some
    irreducible character $\zeta$ of $\Gamma$.
\end{lemma}
\begin{proof}
    \uses{lem:character-inner-product-integer, lem:irr-characters-orthonormal-basis}
    By the accepted basis assumption
    Lemma~\ref{lem:irr-characters-orthonormal-basis}, write $\alpha = \sum_\zeta
    n_\zeta \zeta$ with $n_\zeta \in \mathbb{Z}$ over the irreducible characters
    $\zeta$. Orthonormality (Lemma~\ref{lem:character-inner-product-integer}) gives
    $\langle\alpha,\alpha\rangle_\Gamma = \sum_\zeta n_\zeta^2 = 1$, a sum of
    squares of integers equal to $1$, so exactly one $n_\zeta = \pm 1$ and the
    rest vanish. Thus $\alpha = \pm\zeta$.
\end{proof}

\begin{lemma}[A norm-one virtual character of positive degree is irreducible]
    \label{lem:norm-one-virtual-is-irreducible}
    \lean{LeanEval.GroupTheory.norm_one_virtual_is_irreducible}
    \leanfile{LeanEval/GroupTheory/Frobenius/Character.lean}
    \uses{def:irreducible-character, lem:norm-one-virtual-is-plus-minus-irr, lem:inner-product-eq-dim-hom}
    Let $\alpha$ be a virtual character of a finite group $\Gamma$ (an integer
    combination of irreducible characters) with $\langle\alpha,\alpha\rangle_\Gamma
    = 1$ and $\alpha(1) > 0$. Then $\alpha$ is an irreducible character of
    $\Gamma$, i.e.\ the character of a simple $\mathbb{C}$-representation.
\end{lemma}
\begin{proof}
    \uses{lem:norm-one-virtual-is-plus-minus-irr, lem:inner-product-eq-dim-hom}
    By Lemma~\ref{lem:norm-one-virtual-is-plus-minus-irr}, $\alpha = \pm\zeta$ for
    some irreducible $\zeta$. Since $\alpha(1) > 0$ and $\zeta(1) > 0$ the sign is
    $+$, so $\alpha = \zeta$ is the character of a simple representation.
    (Concretely, a representation $V$ is simple iff
    $\langle\chi_V,\chi_V\rangle_\Gamma = \dim\mathrm{Hom}_\Gamma(V,V) = 1$ by
    Lemma~\ref{lem:inner-product-eq-dim-hom}, i.e.\ Mathlib's
    \texttt{FDRep.simple\_iff\_char\_is\_norm\_one}.)
\end{proof}

\begin{lemma}[The lift is an irreducible character of $G$]
    \label{lem:lift-is-irreducible}
    \lean{LeanEval.GroupTheory.liftClassFun_isIrreducible}
    \leanfile{LeanEval/GroupTheory/Frobenius/Induction.lean}
    \uses{lem:lift-is-virtual-character, lem:lift-norm-one, lem:lift-restricts, lem:norm-one-virtual-is-irreducible, def:irreducible-character}
    For each non-trivial irreducible character $\chi$ of $H$, the lift
    $\widetilde{\chi}$ is an irreducible character of $G$; it is the character of a
    simple complex representation $\widetilde{\rho}$ of $G$ with
    $\widetilde{\chi}(1) = \chi(1)$.
\end{lemma}
\begin{proof}
    \uses{lem:lift-is-virtual-character, lem:lift-norm-one, lem:lift-restricts, lem:norm-one-virtual-is-irreducible}
    The lift is a virtual character (Lemma~\ref{lem:lift-is-virtual-character}) of
    norm one (Lemma~\ref{lem:lift-norm-one}) with positive degree
    $\widetilde{\chi}(1) = \chi(1) > 0$ (Lemma~\ref{lem:lift-restricts}). By
    Lemma~\ref{lem:norm-one-virtual-is-irreducible} it is therefore an irreducible
    character, the character of a simple representation $\widetilde{\rho}$.
\end{proof}

\subsection{The kernel and its normality}

\begin{lemma}[Trace equal to degree forces all eigenvalues to be one]
    \label{lem:char-eq-degree-implies-eigenvalues-one}
    \lean{LeanEval.GroupTheory.char_eq_degree_eigenvalues_one}
    \leanfile{LeanEval/GroupTheory/Frobenius/Character.lean}
    \uses{def:irreducible-character, lem:rep-element-eigenvalues-roots-of-unity}
    With the notation of
    Lemma~\ref{lem:rep-element-eigenvalues-roots-of-unity}, if
    $\chi(g) = \chi(1) = d$ then every eigenvalue $\lambda_i$ of $\rho(g)$ equals
    $1$.
\end{lemma}
\begin{proof}
    \uses{lem:rep-element-eigenvalues-roots-of-unity}
    By Lemma~\ref{lem:rep-element-eigenvalues-roots-of-unity} each
    $|\lambda_i| = 1$, so $d = \sum_i |\lambda_i| \ge |\sum_i \lambda_i| =
    |\chi(g)|$ (\texttt{norm\_sum\_le}). The hypothesis $\chi(g) = d$ makes this an
    equality $|\sum_i \lambda_i| = \sum_i |\lambda_i|$ with all summands of equal
    norm $1$; the equality case of the triangle inequality (the \texttt{SameRay}
    family, e.g.\ \texttt{SameRay.eq\_of\_norm\_eq}) forces all $\lambda_i$ to be
    equal unit vectors. Their common value is then $\chi(g)/d = 1$, so each
    $\lambda_i = 1$.
\end{proof}

\begin{lemma}[Trace equals degree iff the element acts trivially]
    \label{lem:trace-eq-degree-iff-identity}
    \lean{LeanEval.GroupTheory.trace_eq_degree_iff_identity}
    \leanfile{LeanEval/GroupTheory/Frobenius/Character.lean}
    \uses{def:irreducible-character, lem:rep-element-eigenvalues-roots-of-unity, lem:char-eq-degree-implies-eigenvalues-one}
    Let $\rho$ be a finite-dimensional complex representation of a finite group
    $\Gamma$ with character $\chi$. For $g \in \Gamma$, $\rho(g) = 1$ if and only
    if $\chi(g) = \chi(1)$.
\end{lemma}
\begin{proof}
    \uses{lem:rep-element-eigenvalues-roots-of-unity, lem:char-eq-degree-implies-eigenvalues-one}
    If $\rho(g) = 1$ then $\chi(g) = \mathrm{tr}\,1 = \dim\rho = \chi(1)$ by
    \texttt{Representation.char\_one}. Conversely, suppose $\chi(g) = \chi(1) = d$.
    By Lemma~\ref{lem:char-eq-degree-implies-eigenvalues-one} every eigenvalue of
    $\rho(g)$ equals $1$, and $\rho(g)$ is diagonalisable
    (Lemma~\ref{lem:rep-element-eigenvalues-roots-of-unity}); a diagonalisable
    operator all of whose eigenvalues equal $1$ is the identity, so $\rho(g) = 1$.
\end{proof}

\begin{lemma}[Kernel of the lifted homomorphism is a normal subgroup]
    \label{lem:lifted-kernel-is-normal-subgroup}
    \lean{LeanEval.GroupTheory.lifted_kernel_isNormalSubgroup}
    \leanfile{LeanEval/GroupTheory/Frobenius/Induction.lean}
    \uses{lem:lift-is-irreducible}
    For each non-trivial irreducible character $\chi$ of $H$ with lift
    $\widetilde{\rho}$, the set $\{g \in G : \widetilde{\rho}(g) = 1\}$ is a normal
    subgroup of $G$.
\end{lemma}
\begin{proof}
    \uses{lem:lift-is-irreducible}
    By Lemma~\ref{lem:lift-is-irreducible}, $\widetilde{\rho} : G \to
    \mathrm{GL}(V)$ is a group homomorphism. Its kernel $\{g : \widetilde\rho(g) =
    1\}$ is a normal subgroup, namely \texttt{MonoidHom.ker} with
    \texttt{MonoidHom.normal\_ker}.
\end{proof}

\begin{lemma}[Kernel of the lifted representation is normal]
    \label{lem:kernel-normal}
    \lean{LeanEval.GroupTheory.kernel_normal}
    \leanfile{LeanEval/GroupTheory/Frobenius/Induction.lean}
    \uses{lem:lifted-kernel-is-normal-subgroup, lem:trace-eq-degree-iff-identity}
    For each non-trivial irreducible character $\chi$ of $H$, the set
    $\ker \widetilde{\chi} = \{ g \in G : \widetilde{\chi}(g) =
    \widetilde{\chi}(1)\}$ equals $\{g \in G : \widetilde\rho(g) = 1\}$ and is a
    normal subgroup of $G$.
\end{lemma}
\begin{proof}
    \uses{lem:lifted-kernel-is-normal-subgroup, lem:trace-eq-degree-iff-identity}
    By Lemma~\ref{lem:trace-eq-degree-iff-identity}, $\widetilde\chi(g) =
    \widetilde\chi(1) \iff \widetilde\rho(g) = 1$, so the trace-defined set equals
    $\{g : \widetilde\rho(g) = 1\}$, which is a normal subgroup by
    Lemma~\ref{lem:lifted-kernel-is-normal-subgroup}.
\end{proof}

\begin{lemma}[The Frobenius kernel lies in every lifted kernel]
    \label{lem:frobenius-kernel-in-each}
    \lean{LeanEval.GroupTheory.frobenius_kernel_in_each}
    \leanfile{LeanEval/GroupTheory/Frobenius/Induction.lean}
    \uses{def:lifted-class-function, lem:induced-vanishing-restriction, lem:lift-restricts}
    Every $g \in K$ satisfies $\widetilde{\chi}(g) = \widetilde{\chi}(1)$ for each
    irreducible character $\chi$ of $H$; that is, $K \subseteq \ker
    \widetilde{\chi}$.
\end{lemma}
\begin{proof}
    \uses{lem:induced-vanishing-restriction, lem:lift-restricts}
    For $g = 1$ this is trivial. For fixed-point-free $g$,
    Lemma~\ref{lem:induced-vanishing-restriction} gives
    $(\mathrm{Ind}_H^G\theta)(g) = 0$, whence $\widetilde{\chi}(g) = \chi(1) =
    \widetilde{\chi}(1)$ by Lemma~\ref{lem:lift-restricts}.
\end{proof}

\begin{lemma}[Irreducible characters separate non-identity elements from $1$]
    \label{lem:irr-separates}
    \lean{LeanEval.GroupTheory.irr_separates}
    \leanfile{LeanEval/GroupTheory/Frobenius/Induction.lean}
    \uses{def:irreducible-character, lem:irr-characters-orthonormal-basis, lem:trace-eq-degree-iff-identity}
    For $h \in H$ with $h \neq 1$ there is an irreducible character $\chi$ of $H$
    with $\chi(h) \neq \chi(1)$.
\end{lemma}
\begin{proof}
    \uses{lem:irr-characters-orthonormal-basis, lem:trace-eq-degree-iff-identity}
    The regular representation of $H$ is faithful, so $\rho_{\mathrm{reg}}(h) \neq
    1$ for $h \neq 1$. It decomposes into irreducibles (every representation does,
    by Lemma~\ref{lem:irr-characters-orthonormal-basis} / Maschke), so some
    irreducible constituent $\rho$ has $\rho(h) \neq 1$; by
    Lemma~\ref{lem:trace-eq-degree-iff-identity} its character satisfies $\chi(h)
    \neq \chi(1)$. (Equivalently, the common kernel $\bigcap_\chi \ker\chi$ is
    trivial.)
\end{proof}

\begin{lemma}[A point-fixing non-identity element is conjugate into $H \setminus \{1\}$]
    \label{lem:non-kernel-conjugate-to-H}
    \lean{LeanEval.GroupTheory.non_kernel_conjugate_to_H}
    \leanfile{LeanEval/GroupTheory/Frobenius/Basic.lean}
    \uses{lem:partition, def:frobenius-hypotheses, lem:stabilizer-conj}
    Let $g \in G$ with $g \neq 1$ and $g \notin K$. Then $g$ is conjugate in $G$
    to a non-identity element $h \in H$; explicitly there is $s \in G$ with $h =
    s^{-1} g s \in H$ and $h \neq 1$.
\end{lemma}
\begin{proof}
    \uses{lem:partition, lem:stabilizer-conj}
    Since $g \neq 1$ and $g \notin K$, by Lemma~\ref{lem:partition} $g$ fixes
    exactly one point $x \in X$, so $g \in \mathrm{Stab}_G(x)$. By transitivity
    write $x = s \cdot x_0$ for some $s \in G$; then
    Lemma~\ref{lem:stabilizer-conj} gives $\mathrm{Stab}_G(x) = s H s^{-1}$, so
    $h := s^{-1} g s \in H$. Finally $h \neq 1$ because $g \neq 1$ and conjugation
    is injective.
\end{proof}

\begin{lemma}[The intersection of lifted kernels lies in the Frobenius kernel]
    \label{lem:intersection-subset-kernel}
    \lean{LeanEval.GroupTheory.intersection_subset_kernel}
    \leanfile{LeanEval/GroupTheory/Frobenius/Induction.lean}
    \uses{lem:lift-restricts, lem:non-kernel-conjugate-to-H, lem:irr-separates, lem:trace-eq-degree-iff-identity}
    If $g \in \bigcap_{\chi \in \mathrm{Irr}(H)} \ker \widetilde{\chi}$ then $g \in
    K$; that is, $\bigcap_\chi \ker\widetilde\chi \subseteq K$.
\end{lemma}
\begin{proof}
    \uses{lem:lift-restricts, lem:non-kernel-conjugate-to-H, lem:irr-separates, lem:trace-eq-degree-iff-identity}
    Take $g \notin K$ with $g \neq 1$; by
    Lemma~\ref{lem:non-kernel-conjugate-to-H} it is conjugate to a non-identity
    element $h \in H$. By Lemma~\ref{lem:irr-separates} some irreducible $\chi$ of
    $H$ has $\chi(h) \neq \chi(1)$. By Lemma~\ref{lem:lift-restricts} and
    class-invariance of $\widetilde\chi$, $\widetilde\chi(g) = \widetilde\chi(h) =
    \chi(h) \neq \chi(1) = \widetilde\chi(1)$, so $g \notin \ker\widetilde\chi$
    (Lemma~\ref{lem:trace-eq-degree-iff-identity}). Hence $g \notin \bigcap_\chi
    \ker\widetilde\chi$.
\end{proof}

\begin{lemma}[The intersection of lifted kernels is the Frobenius kernel]
    \label{lem:intersection-equals-kernel}
    \lean{LeanEval.GroupTheory.intersection_equals_kernel}
    \leanfile{LeanEval/GroupTheory/Frobenius/Induction.lean}
    \uses{lem:frobenius-kernel-in-each, lem:intersection-subset-kernel}
    The intersection $\bigcap_{\chi \in \mathrm{Irr}(H)} \ker \widetilde{\chi}$,
    over all irreducible characters $\chi$ of $H$, equals $K$ as a subset of $G$.
\end{lemma}
\begin{proof}
    \uses{lem:frobenius-kernel-in-each, lem:intersection-subset-kernel}
    Lemma~\ref{lem:frobenius-kernel-in-each} gives $K \subseteq \bigcap_\chi \ker
    \widetilde{\chi}$, and Lemma~\ref{lem:intersection-subset-kernel} gives the
    reverse inclusion. Hence the two sets are equal.
\end{proof}

\begin{theorem}[Frobenius's theorem: the kernel is normal]
    \label{thm:frobenius-kernel-normal}
    \lean{LeanEval.GroupTheory.frobenius_kernel_isNormal}
    \leanfile{LeanEval/GroupTheory/Frobenius.lean}
    \uses{def:frobenius-kernel, lem:kernel-normal, lem:intersection-equals-kernel}
    There is a normal subgroup $N \trianglelefteq G$ whose underlying set is
    $K = \{1\} \cup \{ g \in G : \forall x \in X,\ g \cdot x \neq x\}$.
\end{theorem}
\begin{proof}
    \uses{lem:kernel-normal, lem:intersection-equals-kernel}
    Let $N = \bigcap_{\chi \in \mathrm{Irr}(H)} \ker \widetilde{\chi}$. Each
    $\ker\widetilde{\chi}$ is a normal subgroup
    (Lemma~\ref{lem:kernel-normal}; for the trivial character $\widetilde\chi =
    \mathbf 1_G$ and $\ker\widetilde\chi = G$), and an indexed intersection of
    normal subgroups is a normal subgroup
    (\texttt{Subgroup.normal\_iInf\_normal}), so $N \trianglelefteq G$. By
    Lemma~\ref{lem:intersection-equals-kernel}, $N = K$ as sets, which is the
    claim.
\end{proof}
